\documentclass[12pt, a4paper]{article}

\usepackage[T2A]{fontenc}
\usepackage[utf8]{inputenc}
\usepackage[bulgarian]{babel}
\usepackage{geometry}
\geometry{a4paper, margin=1in}
\usepackage{tabularx}
\usepackage{xltabular}
\usepackage{ragged2e}
\usepackage{booktabs}
\usepackage{hyperref}

\title{Задание 1: Описание на проекта}
\author{}
\date{}

\begin{document}

\maketitle

\noindent\textbf{Дисциплина:} Проектиране на човеко-машинен интерфейс 2025-2026

\bigskip

\noindent\textbf{Участници в проекта:}
\smallskip

\noindent
\begin{tabularx}{\textwidth}{|l|X|l|X|l|}
\hline
\textbf{№} & \textbf{Име и фамилия} & \textbf{Факултетен №} & \textbf{Специалност} & \textbf{Курс} \\ \hline
1 & Виктор Койчев & 3MI0600164 & Софтуерно инженерство & 3-ти \\ \hline
2 & Емилиан Янев & 0MI0600273 & Софтуерно инженерство & 3-ти \\ \hline
3 & Жан Георгиев & 8MI0600398 & Софтуерно инженерство & 3-ти \\ \hline
4 & Даниел Николов & 0MI0600410 & Софтуерно инженерство & 3-ти \\ \hline
\end{tabularx}

\bigskip

\noindent\textbf{Име на група:}\\
HCI\_2026\_group\_3MI0600164\_0MI0600273\_8MI0600398\_0MI0600410

\bigskip

\noindent\textbf{Име на проекта:}\\
Система за управление и мониторинг на роботизиран склад за фармацевтични продукти (PharmaBot Logistics)

\bigskip
\rule{\textwidth}{0.4pt}

\section*{1. Бизнес нужди и свойства на системата}

\textbf{Контекст и дефиниция на проблема}

Логистиката във фармацевтичния сектор изисква изключителна прецизност, високо ниво на сигурност и непрекъснат температурен контрол (т.нар. „студена верига“). Дори минимални грешки при транспортирането на специализирани медикаменти могат да доведат до сериозни финансови загуби, нарушаване на регулациите и риск за здравето на пациентите. 

В много от съвременните складове процесите все още разчитат на фрагментирани системи, ръчна обработка на документи и визуално проследяване на палетите. Това не само натоварва персонала, но и създава предпоставки за критични човешки грешки. Основният проблем е липсата на единна и интуитивна платформа, която да синхронизира движението на машините с данните за самите пратки.

\bigskip
\noindent\textbf{Предлагано решение}

Проектът представлява специализиран софтуерен интерфейс за управление на флотилия от автономни складови роботи. Целта на системата е да автоматизира процесите по разпределяне, мониторинг и предаване на пратките. Интерфейсът е проектиран с оглед на специфичните нужди и работна среда на три основни групи потребители: мениджър на инвентара, складов оператор и шофьор на камион.

\bigskip
\noindent\textbf{UI/UX акценти и свойства на системата:}
\begin{itemize}
    \item \textbf{Адаптивни интерфейси според ролята:} Десктоп приложение с възможност за тъмна тема (за намаляване умората на очите) за нуждите на операторите в контролната зала, и мобилно приложение с висок контраст и големи touch-елементи за шофьорите на рампата, съобразено с работа с ръкавици.
    \item \textbf{Визуален мониторинг:} Интерактивна 2D карта, която предоставя бърз достъп до статуса на роботите чрез цветови индикатори (зелено за нормално движение, жълто при зареждане, червено за възникнал проблем).
    \item \textbf{Geo-fencing управление:} Възможност за ръчно очертаване (drag-and-drop) на рискови зони върху картата при възникване на инциденти (напр. разлив). Системата автоматично генерира нови, безопасни маршрути за роботите.
    \item \textbf{Контрол на „студената верига“:} При засичане на температурно отклонение, системата блокира екрана с критично предупреждение, което изисква изрично действие и потвърждение от оператор.
    \item \textbf{Безхартиена логистика:} Идентификацията и предаването на товара се извършват изцяло дигитално чрез сканиране на QR код от дисплея на робота и прикачване на снимка през мобилното устройство на шофьора.
\end{itemize}

\bigskip
\noindent\textbf{Основни ползи за участниците в процеса:}

\begin{itemize}
    \item \textit{За мениджърския състав:} Осигурява се пълна проследяемост на всяко действие, минимизират се загубите от развален инвентар благодарение на температурния контрол и се повишава сигурността чрез дигитални доказателства за доставка.
    \item \textit{За складовите оператори:} Намалява се когнитивното натоварване чрез визуализиране на данните върху карта вместо в сложни таблици. Инструментите за директна намеса позволяват бързо овладяване на аварийни ситуации.
    \item \textit{За шофьорите:} Процесът по приемане на стоката е сведен до две прости стъпки (сканиране и снимане), без необходимост от изчакване и попълване на физически документи.
\end{itemize}

\rule{\textwidth}{0.4pt}

\section*{2. Кратко описание на потребителските случаи (Use cases)}

\noindent
\begin{xltabular}{\textwidth}{|>{\RaggedRight\arraybackslash}p{4cm}|>{\RaggedRight\arraybackslash}X|>{\RaggedRight\arraybackslash}X|}
\hline
\textbf{Потребителски случай} & \textbf{Описание на процеса} & \textbf{Актьори и роли} \\ \hline
\endfirsthead

\hline
\textbf{Потребителски случай} & \textbf{Описание на процеса} & \textbf{Актьори и роли} \\ \hline
\endhead

\textbf{2.1. Вход в системата с 2FA} & Мениджърът въвежда своите идентификационни данни и код за сигурност. След успешна верификация, системата предоставя достъп до контролния панел за скъпоструващи медикаменти. & \textbf{Мениджър инвентар} – контролира достъпа и отговаря за сигурността на продуктите. \\ \hline

\textbf{2.2. Възлагане на задача на робот} & Операторът или мениджърът селектира поръчка за изпълнение. Системата анализира локациите и възлага задачата на оптималния свободен робот с достатъчен заряд на батерията. & \textbf{Мениджър инвентар} – дефинира приоритети.\par\medskip \textbf{Системата} – изчислява маршрут и ресурс. \\ \hline

\textbf{2.3. Наблюдение на склада в реално време} & Операторът зарежда интерактивната карта на склада, където може да мащабира (zoom) изгледа и да следи движението и състоянието на активните роботи. & \textbf{Складов оператор} – осъществява визуален контрол върху логистиката. \\ \hline

\textbf{2.4. Преглед на детайли за конкретен робот} & При клик върху иконата на даден робот от картата, се отваря страничен панел със справка за ниво на батерията, тип на товара, текуща температура и крайна цел. & \textbf{Складов оператор} – следи експлоатационните параметри на машините. \\ \hline

\textbf{2.5. Създаване на динамична зона за опасност} & При инцидент в склада, операторът очертава блокирания участък върху картата. Системата регистрира зоната и незабавно пренасочва трафика. & \textbf{Складов оператор} – предотвратява задръствания и аварии в обекта. \\ \hline

\textbf{2.6. Обработка на температурна аларма} & Системата засича температурно отклонение и блокира екрана с аларма. Операторът потвърждава събитието и пренасочва компрометираната пратка към карантинна зона. & \textbf{Складов оператор} – гарантира спазването на изискванията на "студената верига". \\ \hline

\textbf{2.7. Вход в мобилното приложение} & Шофьорът отваря приложението на служебния таблет или телефон, за да получи достъп до списъка с очаквани палети за текущия курс. & \textbf{Шофьор на камион} – организира процеса по натоварване. \\ \hline

\textbf{2.8. Идентификация чрез сканиране на QR код} & За да удостовери, че е легитимният получател, шофьорът използва камерата на устройството си, за да сканира уникалния QR код, генериран на дисплея на робота. & \textbf{Шофьор на камион} – инициатор на предаването.\par\medskip \textbf{Робот} – верифицира самоличността. \\ \hline

\textbf{2.9. Създаване на дигитално доказателство} & След физическото прехвърляне на пратката в камиона, шофьорът прави снимка през приложението, която служи като електронен протокол за приемане. & \textbf{Шофьор на камион} – документира завършената операция. \\ \hline

\textbf{2.10. Финализиране на експедицията} & След получаване на сканирания код и снимката, системата променя статуса на поръчката на "Изпълнена" и освобождава робота за следващи задачи. & \textbf{Системата} – обработва данните от транзакцията.\par\medskip \textbf{Мениджър инвентар} – получава нотификация. \\ \hline

\end{xltabular}

\end{document}